% GENERATED FILE. Edit canonical lessons, then run npm run generate:books.
% canonical-source-hash: fnv1a64:b0535981b68b3a06
% canonical-lessons: JA-C496-yawarakai, JA-C496-shiokarai, JA-C496-nama, JA-C496-enbun, JA-C496-shokuzai

\chapter{Soft and Salty}
\label{ch:ja-soft-and-salty}
\hlchaptermodality{496}

\begin{chapteropening}
\textbf{By the end of this chapter:} \emph{I can say what food is like and what it contains.}

Complete the last lesson of chapter 496: I can say what food is like and what it contains.
\end{chapteropening}

\section[\texorpdfstring{\ja{やわらかい} (yawarakai)}{やわらかい (yawarakai)}]{\ja{やわらかい} (yawarakai) — soft; tender}

\label{lesson:JA-C496-yawarakai}



\begin{quote}
Before the new one: say the Japanese for fresh, then the Japanese for light (of food).
\end{quote}

\subsection*{You'll want to know: \ja{やわらかい}}
\textbf{\ja{やわらかい}} — \emph{yawarakai} — \textquotedblleft{}soft; tender\textquotedblright{}.

\begin{quote}
\textbf{\ja{この} \ja{にくは} \ja{やわらかいです。}} — \emph{kono niku wa yawarakai desu} — \textquotedblleft{}This meat is tender.\textquotedblright{}
\end{quote}

An i-adjective.

\begin{cousinweb}
\textbf{\ja{やわらかい}} (\emph{yawarakai}) is a native Japanese word (wago), inherited rather than borrowed from Chinese; its kanji was attached later as a kun'yomi, and its older origin is not settled (JMdict, KANJIDIC2).
\end{cousinweb}

\begin{grammarlens}[title={What you've built}]
One more word for describing what you lived and felt.
\end{grammarlens}

\subsection*{Your turn}
\begin{itemize}
\raggedright
  \item \emph{Say it:} \emph{yawarakai}
  \item \emph{Say it:} \emph{yawarakai}, once more
  \item \emph{Say it:} \emph{kono niku wa yawarakai desu} — This meat is tender.
  \item \emph{Recall:} say the Japanese for fresh, then the Japanese for light (of food), then say \emph{yawarakai} again
\end{itemize}

\begin{tcolorbox}[breakable,colback=teal!4,colframe=teal!35!black,title={Before you move on}]
What does \ja{やわらかい} mean? (\textquotedblleft{}Soft; tender\textquotedblright{}.) Say it once more.
\end{tcolorbox}
\section[\texorpdfstring{\ja{しおからい} (shiokarai)}{しおからい (shiokarai)}]{\ja{しおからい} (shiokarai) — salty}

\label{lesson:JA-C496-shiokarai}



\begin{quote}
Before the new one: say the Japanese for light (of food), then the Japanese for soft.
\end{quote}

\subsection*{You'll want to know: \ja{しおからい}}
\textbf{\ja{しおからい}} — \emph{shiokarai} — \textquotedblleft{}salty\textquotedblright{}.

\begin{quote}
\textbf{\ja{この} \ja{さかなは} \ja{しおからいです。}} — \emph{kono sakana wa shiokarai desu} — \textquotedblleft{}This fish is salty.\textquotedblright{}
\end{quote}

An i-adjective.

\begin{cousinweb}
\textbf{\ja{しおからい}} (\emph{shiokarai}) is a native Japanese compound (wago), inherited rather than borrowed from Chinese: built from \emph{shio}, \textquotedblleft{}salt\textquotedblright{} and \emph{karai}, \textquotedblleft{}hot, spicy\textquotedblright{} (JMdict, KANJIDIC2).
\end{cousinweb}

\begin{grammarlens}[title={What you've built}]
One more word for describing what you lived and felt.
\end{grammarlens}

\subsection*{Your turn}
\begin{itemize}
\raggedright
  \item \emph{Say it:} \emph{shiokarai}
  \item \emph{Say it:} \emph{shiokarai}, once more
  \item \emph{Say it:} \emph{kono sakana wa shiokarai desu} — This fish is salty.
  \item \emph{Recall:} say the Japanese for light (of food), then the Japanese for soft, then say \emph{shiokarai} again
\end{itemize}

\begin{tcolorbox}[breakable,colback=teal!4,colframe=teal!35!black,title={Before you move on}]
What does \ja{しおからい} mean? (\textquotedblleft{}Salty\textquotedblright{}.) Say it once more.
\end{tcolorbox}
\section[\texorpdfstring{\ja{なま} (nama)}{なま (nama)}]{\ja{なま} (nama) — raw; uncooked}

\label{lesson:JA-C496-nama}



\begin{quote}
Before the new one: say the Japanese for soft, then the Japanese for salty.
\end{quote}

\subsection*{You'll want to know: \ja{なま}}
\textbf{\ja{なま}} — \emph{nama} — \textquotedblleft{}raw; uncooked\textquotedblright{}.

\begin{quote}
\textbf{\ja{にほんで} \ja{なまの} \ja{たまごを} \ja{たべます。}} — \emph{nihon de nama no tamago o tabemasu} — \textquotedblleft{}In Japan people eat raw eggs.\textquotedblright{}
\end{quote}

\begin{cousinweb}
\textbf{\ja{なま}} (\emph{nama}) is a native Japanese word (wago), inherited rather than borrowed from Chinese; its kanji was attached later as a kun'yomi, and its older origin is not settled (JMdict, KANJIDIC2).
\end{cousinweb}

\begin{grammarlens}[title={What you've built}]
One more word for describing what you lived and felt.
\end{grammarlens}

\subsection*{Your turn}
\begin{itemize}
\raggedright
  \item \emph{Say it:} \emph{nama}
  \item \emph{Say it:} \emph{nama}, once more
  \item \emph{Say it:} \emph{nihon de nama no tamago o tabemasu} — In Japan people eat raw eggs.
  \item \emph{Recall:} say the Japanese for soft, then the Japanese for salty, then say \emph{nama} again
\end{itemize}

\begin{tcolorbox}[breakable,colback=teal!4,colframe=teal!35!black,title={Before you move on}]
What does \ja{なま} mean? (\textquotedblleft{}Raw; uncooked\textquotedblright{}.) Say it once more.
\end{tcolorbox}
\section[\texorpdfstring{\ja{えんぶん} (enbun)}{えんぶん (enbun)}]{\ja{えんぶん} (enbun) — salt content}

\label{lesson:JA-C496-enbun}



\begin{quote}
Before the new one: say the Japanese for salty, then the Japanese for raw.
\end{quote}

\subsection*{You'll want to know: \ja{えんぶん}}
\textbf{\ja{えんぶん}} — \emph{enbun} — \textquotedblleft{}salt content\textquotedblright{}.

\begin{quote}
\textbf{\ja{この} \ja{みそしるは} \ja{えんぶんが} \ja{おおいです。}} — \emph{kono misoshiru wa enbun ga ōi desu} — \textquotedblleft{}This miso soup has a lot of salt in it.\textquotedblright{}
\end{quote}

\begin{cousinweb}
\textbf{\ja{えんぶん}} (\emph{enbun}) is Sino-Japanese (kango): two Chinese characters read with their on'yomi, readings borrowed from Chinese. They are \emph{en}, \textquotedblleft{}salt\textquotedblright{} (Middle Chinese \emph{yem}) and \emph{bun}, \textquotedblleft{}part\textquotedblright{} (\emph{bjunH}) (KANJIDIC2; Middle Chinese in Baxter's transcription).
\end{cousinweb}

\begin{grammarlens}[title={What you've built}]
One more word for describing what you lived and felt.
\end{grammarlens}

\subsection*{Your turn}
\begin{itemize}
\raggedright
  \item \emph{Say it:} \emph{enbun}
  \item \emph{Say it:} \emph{enbun}, once more
  \item \emph{Say it:} \emph{kono misoshiru wa enbun ga ōi desu} — This miso soup has a lot of salt in it.
  \item \emph{Recall:} say the Japanese for salty, then the Japanese for raw, then say \emph{enbun} again
\end{itemize}

\begin{tcolorbox}[breakable,colback=teal!4,colframe=teal!35!black,title={Before you move on}]
What does \ja{えんぶん} mean? (\textquotedblleft{}Salt content\textquotedblright{}.) Say it once more.
\end{tcolorbox}
\section[\texorpdfstring{\ja{しょくざい} (shokuzai)}{しょくざい (shokuzai)}]{\ja{しょくざい} (shokuzai) — an ingredient; foodstuff}

\label{lesson:JA-C496-shokuzai}



\begin{quote}
Before the new one: say the Japanese for raw, then the Japanese for salt content.
\end{quote}

\subsection*{You'll want to know: \ja{しょくざい}}
\textbf{\ja{しょくざい}} — \emph{shokuzai} — \textquotedblleft{}an ingredient; foodstuff\textquotedblright{}.

\begin{quote}
\textbf{\ja{いちばで} \ja{しょくざいを} \ja{かいます。}} — \emph{ichiba de shokuzai o kaimasu} — \textquotedblleft{}I buy ingredients at the market.\textquotedblright{}
\end{quote}

\begin{cousinweb}
\textbf{\ja{しょくざい}} (\emph{shokuzai}) is Sino-Japanese (kango): two Chinese characters read with their on'yomi, readings borrowed from Chinese. They are \emph{shoku}, \textquotedblleft{}eat\textquotedblright{} (Middle Chinese \emph{zyik}) and \emph{zai}, \textquotedblleft{}ingredients\textquotedblright{} (KANJIDIC2; Middle Chinese in Baxter's transcription).
\end{cousinweb}

\begin{grammarlens}[title={What you've built}]
One more word for describing what you lived and felt.
\end{grammarlens}

\subsection*{Your turn}
\begin{itemize}
\raggedright
  \item \emph{Say it:} \emph{shokuzai}
  \item \emph{Say it:} \emph{shokuzai}, once more
  \item \emph{Say it:} \emph{ichiba de shokuzai o kaimasu} — I buy ingredients at the market.
  \item \emph{Recall:} say the Japanese for raw, then the Japanese for salt content, then say \emph{shokuzai} again
\end{itemize}

\begin{tcolorbox}[breakable,colback=teal!4,colframe=teal!35!black,title={Before you move on}]
What does \ja{しょくざい} mean? (\textquotedblleft{}An ingredient; foodstuff\textquotedblright{}.) Say it once more.
\end{tcolorbox}
